\exercisesection{}

\begin{enumerate}
\def\labelenumi{\arabic{enumi})}
\tightlist
\item
  设 A 是 \(\mathbf{C}\) 的紧子集，它等于其内部 V 的闭包。我们置 \(\mathcal{C}(\mathrm{A}) = \mathcal{C}(\mathrm{A};\mathbf{C})\) 。设 \(B(A)\) 是 \(\mathcal{C}(\mathrm{A})\) 中由在 V 上的限制为全纯的函数 \(f\in \mathcal{C}(\mathrm{A})\) 所组成的向量子空间；
\end{enumerate}

它是 \(\mathcal{C}(\mathrm{A})\) 的一个巴拿赫子代数。设 \(u \in \mathcal{L}(\mathrm{B}(\mathrm{A}))\) 是满足 \(u(f)(z) = zf(z)\) （对所有 \(z \in \mathrm{A}\) 及每个 \(f \in \mathrm{B}(\mathrm{A})\) ）的线性映射。

\begin{enumerate}
\def\labelenumi{\alph{enumi})}
\item
  \(u\) 的谱等于 A。
\item
  对每个 \(z \in V\) ，B(A) 的自同态 u - z 是指标为 1 的弗雷德霍姆自同态。
\item
  对每个 \(z \in A - V\) ，B(A) 的自同态 u - z 不是弗雷德霍姆自同态。
\item
  u 的稳定谱是 A 的边界，u 的敏感谱是空集，而 u 的本质谱是 A。
\item
  我们假设 A 在 \(\mathbf{C}\) 中的余集是连通的。对每个 \(h \in \mathcal{L}^{c}(\mathrm{B}(\mathrm{A}))\) ，\(u + h\) 的本质谱等于 A。
\end{enumerate}

\begin{enumerate}
\def\labelenumi{\arabic{enumi})}
\setcounter{enumi}{1}
\tightlist
\item
  我们保留上一习题的记号。假设 A 包含单位圆 \(U\) （\(\mathbf{C}\) 中的），且 \(0 \notin A\) 。用 \(\varphi_{A}\) 表示 A 的特征函数。
\end{enumerate}

\begin{enumerate}
\def\labelenumi{\alph{enumi})}
\tightlist
\item
  映射
\end{enumerate}

\[
\ell \colon f \mapsto \frac {1}{2 i \pi} \int_ {\mathbf {U}} f (z) d z
\]

是 B(A) 上的一个连续线性形式。

\begin{enumerate}
\def\labelenumi{\alph{enumi})}
\setcounter{enumi}{1}
\tightlist
\item
  设 \(h\) 是 B(A) 的有限秩自同态，满足 \(h(f) = \ell(f)\varphi_{\mathrm{A}}\) 。\(u - h\) 的本质谱异于 \(u\) 的本质谱。（证明它包含 0。）
\end{enumerate}

\begin{enumerate}
\def\labelenumi{\arabic{enumi})}
\setcounter{enumi}{2}
\tightlist
\item
  设 E 是复希尔伯特空间。设 U 是 \(\mathbf{C}\) 的连通开子集。设 f 是从 U 到空间 \(\mathcal{L}^{c}(\mathrm{E})\) 中的全纯函数。假设存在 \(z \in U\) 使得 1 属于 \(f(z)\) 的预解集。
\end{enumerate}

\begin{enumerate}
\def\labelenumi{\alph{enumi})}
\item
  由 \(z \in U\) （满足 \(1 \in \operatorname{Sp}(f(z))\) ）组成的集 D 是 U 的一个离散子集。
\item
  从 U-D 到 \(\mathcal{L}(\mathrm{E})\) 中的由 \(z\mapsto \mathrm{R}(f(z),1)\) 定义的映射是全纯的。
\end{enumerate}

我们固定一点 \(z_0 \in \mathrm{D}\) 。

\begin{enumerate}
\def\labelenumi{\alph{enumi})}
\setcounter{enumi}{2}
\item
  存在整数 \(k \geqslant 1\) ，使得从 U-D 到 \(\mathcal{L}(\mathrm{E})\) 中的由 \(z \mapsto (z - z_{0})^{k}\mathrm{R}(f(z), 1)\) 定义的映射可延拓为 \((\mathrm{U}-\mathrm{D}) \cup \{z_{0}\}\) 上的一个全纯映射。
\item
  存在 E 的自同态的唯一序列 \((u_{n})_{n\in\mathbf{N}}\) （当 n 充分大时 \(u_{n}=0\) ），以及从 \((\mathrm{U}-\mathrm{D})\cup\{z_{0}\}\) 到 \(\mathcal{L}(\mathrm{E})\) 中的唯一全纯函数 S，使得对所有 \(z\in U-D\) 有
\end{enumerate}

\[
\mathrm{R} (f (z), 1) = \sum_ {j = - k} ^ {- 1} (z - z _ {0}) ^ {j} u _ {j} + \mathrm{S} (z).
\]

\begin{enumerate}
\def\labelenumi{\alph{enumi})}
\setcounter{enumi}{4}
\tightlist
\item
  自同态 \(u_{-1}\) 是 \(f(z_0)\) 相对于 \(\{1\}\) （\(f(z_0)\) 的谱的开且闭子集）的谱投影。
\end{enumerate}

\begin{enumerate}
\def\labelenumi{\arabic{enumi})}
\setcounter{enumi}{3}
\tightlist
\item
  设 E 是可数型的巴拿赫空间。设 \(u \in \mathcal{L}(\mathrm{E})\) 。假设序列 \((\| u^k \|)_{k \in \mathbf{N}}\) 有界。
\end{enumerate}

对每个 \(\lambda\in\mathrm{Sp}(u)\cap\mathbf{U}\) ，设 \(x_{\lambda}\in\mathrm{E}\) 是 \(u\) 的属于 \(\lambda\) 的范数为 1 的特征向量 。

\begin{enumerate}
\def\labelenumi{\alph{enumi})}
\tightlist
\item
  对所有 \(\lambda_1\) 、\(\lambda_2\) （属于 \(\mathrm{Sp}(u) \cap \mathbf{U}\) ）以及每个整数 \(k \in \mathbf{N}\) ，我们有
\end{enumerate}

\[
\left| \lambda_ {1} ^ {k} - \lambda_ {2} ^ {k} \right| \leqslant \left(1 + \| u ^ {k} \|\right) \| x _ {\lambda_ {1}} - x _ {\lambda_ {2}} \|.
\]

\begin{enumerate}
\def\labelenumi{\alph{enumi})}
\setcounter{enumi}{1}
\item
  存在实数 c \textgreater{} 0 ，使得 \(\|x_{\lambda_{1}} - x_{\lambda_{2}}\| \geqslant c\) 对所有 \(\lambda_{1} \neq \lambda_{2}\) （\(\operatorname{Sp}(u) \cap \mathbf{U}\) 中的）成立。
\item
  集合 \(\operatorname{Sp}(u) \cap \mathbf{U}\) 是可数的（“贾米森定理”）。
\item
  若 E 不是可数型的，贾米森定理一般不成立。
\end{enumerate}

\begin{enumerate}
\def\labelenumi{\arabic{enumi})}
\setcounter{enumi}{4}
\tightlist
\item
  本习题的目标是在空间 E 为巴拿赫空间的情形给出洛莫诺索夫定理（III, p. 13 的推论 2）的另一个证明。
\end{enumerate}

设 E 是维数至少为 2 的复巴拿赫空间。设 h 是 E 的一个非零紧自同态，A 是 h 在 \(\mathcal{L}(\mathrm{E})\) 中的交换子。a) 下列条件之一成立：

\begin{enumerate}
\def\labelenumi{(\roman{enumi})}
\item
  存在 \(y \in E\) ，使得 Ay 是 E 中的一个非零且不稠密的子空间；
\item
  对 A 的每个内部非空的闭有界子集 C，存在 A 的有限子集 I，使得 \(h(\mathrm{C})\) 包含于诸集合 \(u^{-1}(\mathrm{C})\) （\(u \in I\) ）的并之中。
\end{enumerate}

\begin{enumerate}
\def\labelenumi{\alph{enumi})}
\setcounter{enumi}{1}
\tightlist
\item
  假设上一问题的条件 (ii) 成立。设 C 是 E 的一个内部非空的闭有界子集。对每个整数 \(m \geqslant 1\) ，存在 A 中的有限序列 \((u_{1}, \ldots, u_{m})\) 使得
\end{enumerate}

\[
u _ {1} \circ \dots \circ u _ {m} \circ h ^ {m} (x _ {0}) \in \mathrm{C}.
\]

（打乒乓球。）

\begin{enumerate}
\def\labelenumi{\alph{enumi})}
\setcounter{enumi}{2}
\item
  在同一假设下，证明 h 的谱半径非零（在相反的情形，考虑一个元素 \(x_{0} \in E\) ，满足 \(\|h(x_{0})\| > \|h\|\) ，以及以 \(x_{0}\) 为心、\(\|h\|\) 为半径的闭球 C；注意到 \(\|( \lambda h)^{n}\|\) 对每个 \(\lambda \in R\) 趋于 0，证明上一问题的结论将蕴含 0 属于 C 的闭包）。
\item
  对 E 的每个与 E 的某个非零紧自同态交换的自同态 \(u\) ，存在异于 E 的非零闭子空间 \(F \subset E\) 使得 \(u(F) \subset F\) 。
\end{enumerate}

\begin{enumerate}
\def\labelenumi{\arabic{enumi})}
\setcounter{enumi}{5}
\tightlist
\item
  设 A 是一个幺 Z-代数。假设 A 是自由 Z-模。设 B 是 A 的一个子集，它作为 Z-模是 A 的基。我们称对子 \((A, B)\) 是一个自然代数，如果 A 关于 B 的结构常数（A, III, p. 10）属于 N，且 \(1_{A}\) 在基 B 中的系数属于 N。此外，我们称 (A, B) 是幺的，如果 \(1_{A}\) 是 B 的元素。
\end{enumerate}

设 (A, B) 是一个自然代数。设 \(B_{0}\) 是由 \(b \in B\) （使 b 在 \(1_{A}\) 中的系数非零）组成的集。用 \(\tau_{A}\) 表示从 A 到 Z 的唯一的 Z-模同态，使得 \(\tau_{\mathrm{A}}(b) = 1\) （若 \(b \in B_{0}\) ），否则 \(\tau_{\mathrm{A}}(b) = 0\) 。

\begin{enumerate}
\def\labelenumi{\alph{enumi})}
\item
  \(\tau_{\mathrm{A}}(xy) = \tau_{\mathrm{A}}(yx)\) 对所有 \((x,y)\in \mathrm{A}^2\) 成立。
\item
  对所有 \(b \neq b'\) （在 \(B_{0}\) 中），有 \(b^2 = b\) 且 \(bb' = 0\) 。
\end{enumerate}

\begin{enumerate}
\def\labelenumi{\arabic{enumi})}
\setcounter{enumi}{6}
\tightlist
\item
  设 (A, B) 是一个自然代数。
\end{enumerate}

\begin{enumerate}
\def\labelenumi{\alph{enumi})}
\tightlist
\item
  至多存在一个对合 \(\iota\) （\(\mathcal{B}\) 的），使得从 A 到 A 的唯一 \(\mathbf{Z}\)-模同态（把 \(b\) 映为 \(\iota(b)\) ）是 A 的一个反自同构，且满足条件 \(\tau_{\mathrm{A}}(bb') = 1\) （若 \(b' = \iota(b)\) ），否则 \(\tau_{\mathrm{A}}(bb') = 0\) 。
\end{enumerate}

当这样一个对合 \(\iota\) 存在时，我们称 (A, \(\mathcal{B}\) ) 是一个自然基础代数。

我们现在假设 (A, B) 是一个自然基础代数。

\begin{enumerate}
\def\labelenumi{\alph{enumi})}
\setcounter{enumi}{1}
\tightlist
\item
  于是 \(\iota(b) = b\) 对所有 \(b \in B_{0}\) 成立，并有公式
\end{enumerate}

\[
1 _ {\mathrm{A}} = \sum_ {b \in \mathcal {B} _ {0}} b.
\]

\begin{enumerate}
\def\labelenumi{\alph{enumi})}
\setcounter{enumi}{2}
\tightlist
\item
  A 的结构常数 \(\gamma_{bb'}^{b''}\) 满足 \(\gamma_{bb'}^{\iota(b'')} = \tau_{\mathrm{A}}(bb'b'')\) （对 \((b, b', b'')\) ，\(\mathcal{B}^{3}\) 中的）。
\end{enumerate}

\begin{enumerate}
\def\labelenumi{\arabic{enumi})}
\setcounter{enumi}{7}
\tightlist
\item
  设 (A, B) 是一个自然基础代数，且 A 的秩有限。
\end{enumerate}

\begin{enumerate}
\def\labelenumi{\alph{enumi})}
\tightlist
\item
  对每个 \(x \in A\) ，元素
\end{enumerate}

\[
\sum_ {b \in \mathcal {B}} b x \iota (b)
\]

是中心元。

在本习题的其余部分，我们进一步假设 (A, \(\mathcal{B}\) ) 是幺的。

\begin{enumerate}
\def\labelenumi{\alph{enumi})}
\setcounter{enumi}{1}
\item
  对所有 \(b\) 与 \(b'\) （\(\mathcal{B}\) 中的），存在 \(b_1\) 与 \(b_2\) （\(\mathcal{B}\) 中的），使得 \(b'\) 在 \(bb_1\) 中以及在 \(b_2b\) 中的系数非零。
\item
  对每个 \(x \in B\) ，在基 B 中表示从 A 到 A 的映射 \(y \mapsto xy\) 的矩阵是一个实矩阵，其所有系数都是正的。用 pf 表示从 A 到 R 的唯一同态，使得 \(\operatorname{pf}(b)\) （当 \(b \in B\) 时）是这个矩阵的谱半径。
\end{enumerate}

\(d)\) 对每个 \(b \in \mathcal{B}\) ，实数 \(\alpha = \mathrm{pf}(b)\) 是一个代数整数，满足 \(\alpha \geqslant 1\) ；对每个共轭数 \(\alpha' \in \mathbf{C}\) （\(\alpha\) 的），有 \(|\alpha'| \leqslant \alpha\) 。

\begin{enumerate}
\def\labelenumi{\alph{enumi})}
\setcounter{enumi}{4}
\tightlist
\item
  设 \(a_0 \in \mathrm{A}\) 是 \(\mathcal{B}\) 的诸元素之和。表示映射 \(f: y \mapsto ya_0\) 的矩阵在基 \(\mathcal{B}\) 中的所有系数都严格为正。存在一个元素 \(r \in \mathbf{C} \otimes \mathrm{A}\) ，它在基 \((1 \otimes b)_{b \in \mathcal{B}}\) 中的系数都严格为正，使得 \(f_{\mathbf{C}}(r) = \varrho(f_{\mathbf{C}})r\) 。
\end{enumerate}

\(f)\) 对每个 \(x \in \mathrm{A}\) ，存在唯一的复数 \(\delta(x)\) ，使得 \(xr = \delta(x)r\) 在 \(\mathbf{C} \otimes \mathrm{A}\) 中成立。映射 \(\delta: \mathrm{A} \to \mathbf{C}\) 是 \(\mathrm{A}\) 的一个特征标。

\begin{enumerate}
\def\labelenumi{\alph{enumi})}
\setcounter{enumi}{6}
\tightlist
\item
  \(\delta(x) = \mathrm{pf}(x)\) 对所有 \(x \in \mathrm{A}\) 成立。特别地，映射 \(x \mapsto \mathrm{pf}(x)\) 是 A 的一个特征标。
\end{enumerate}

\(h)\) 映射 pf 是 A 的唯一特征标，使得 \(\mathrm{pf}(b) \geqslant 0\) 对所有 \(b \in \mathcal{B}\) 成立；此外，\(\mathrm{pf}(b) > 0\) 对所有 \(b \in \mathcal{B}\) 成立。

\begin{enumerate}
\def\labelenumi{\roman{enumi})}
\tightlist
\item
  对每个 \(y \in \mathrm{A}\) ，有 \(ry = \delta(y)r\) （在 \(\mathbf{C} \otimes \mathbf{A}\) 中）。
\end{enumerate}

\begin{enumerate}
\def\labelenumi{\alph{enumi})}
\setcounter{enumi}{9}
\tightlist
\item
  元素
\end{enumerate}

\[
r _ {0} = \sum_ {b \in \mathcal {B}} \mathrm{pf} (b) b
\]

具有问题 i) 所要求的性质。

\begin{enumerate}
\def\labelenumi{\arabic{enumi})}
\setcounter{enumi}{8}
\tightlist
\item
  \begin{enumerate}
  \def\labelenumii{\alph{enumii})}
  \tightlist
  \item
    设 \(n\) 是整数，\(M\) 是系数属于 \(\mathbf{N}\) 的一个方阵。假设 \(\varrho(M^t M) = \varrho(M)^2\) 且 \(|\varrho(M)| < 2\) 。则存在整数 \(m \geqslant 2\) 使得 \(\varrho(M) = 2\cos(\pi/m)\) 。（利用克罗内克定理，参见 A, V, p. 155, 习题 17。）
  \end{enumerate}
\end{enumerate}

\begin{enumerate}
\def\labelenumi{\alph{enumi})}
\setcounter{enumi}{1}
\item
  对每个整数 \(m \geqslant 2\) ，存在一个系数属于 N 的方阵 M，使得 \(\varrho(M) = 2 \cos(\pi/m)\) 。
\item
  设 (A, B) 是一个秩有限的幺自然基础代数。设 \(b \in B\) 。若 \(\operatorname{pf}(b) < 2\) ，则存在整数 \(m \geqslant 3\) 使得 \(\operatorname{pf}(b) = 2 \cos(\pi/m)\) 。
\end{enumerate}

习题 10 至 24 给出 III, p. 100 的推论 5 的另一个证明以及若干强化。其中将用到下列概念。

设 I 是有限集，n 是它的基数。

我们用 \(1_{\mathrm{I}}\) 表示 (I, I) 型的实单位矩阵。

- 一个 (I, I) 型的实系数方阵称为正的，如果它的所有系数都 \(\geqslant 0\) ；称为严格正的，如果它的所有系数都 \(>0\) 。

- (I,I) 型的实系数方阵 \((a_{ij})_{i,j\in \mathrm{I}}\) 称为可约的，如果存在把 I 分成两个非空子集 I' 与 I'\,' 的划分，使得 \(a_{ij} = 0\) 对所有 \((i,j)\in \mathrm{I}'\times \mathrm{I}''\) 成立；若它不是可约的，则称为不可约的。

- 我们把 \(\mathbf{R}^{\mathrm{I}}\) 的由典范基向量的某个子集生成的任一向量子空间称为 \(\mathbf{R}^{\mathrm{I}}\) 的坐标子空间。坐标子空间共有 \(2^{n}\) 个；对每个整数 \(k\in\{0,\ldots,n\}\) ，其中维数为 \(k\) 的坐标子空间有 \(\binom{n}{k}\) 个。

因此，(I, I) 型的实矩阵是不可约的，当且仅当它所稳定的坐标子空间只有 \(\{0\}\) 和 \(R^{I}\) 。

对所有 \(x\) 与 \(y\) （\(\mathbf{R}^{\mathrm{I}}\) 中的），我们记 \(x \preccurlyeq y\) （及 \(y \succcurlyeq x\) ），如果 \(y - x\) 属于 \(\mathbf{R}^{\mathrm{I}}\) 中坐标为正的元素所成的锥，亦即 \(x_i \leqslant y_i\) 对所有 \(i \in \mathrm{I}\) 成立。我们记 \(x \prec y\) （及 \(y \succ x\) ），如果 \(y - x\) 的所有坐标都严格为正。

\begin{enumerate}
\def\labelenumi{\arabic{enumi})}
\setcounter{enumi}{9}
\item
  设 \(A\) 是 (I, I) 型的正且不可约的实矩阵。对每个整数 \(n\) （满足 \(n \geqslant \operatorname{Card}(\mathrm{I})\) ），矩阵 \((1_{\mathrm{I}} + A)^{n-1}\) 是严格正的。（设 \(y \in \mathbf{R}_{+}^{n}\) ，\(z = (1_{\mathrm{I}} + A)y\) ；若 \(y \neq 0\)，证明由 \(i \in \mathrm{I}\) （使 \(z_{i} = 0\) ）所成的集严格包含于 \(\{i \in \mathrm{I} \mid y_{i} = 0\}\) 。）
\item
  设 \(A\) 是 (I, I) 型的正且不可约的实矩阵。用 \(\psi \in \mathbf{R}[\mathrm{X}]\) 表示 \(A\) 的极小多项式，并置 \(m = \deg(\psi)\) 。
\end{enumerate}

对每个整数 \(q\)，置 \(A^q = (a_{ij}^{(q)})_{i,j \in \mathrm{I}}\) 。

对每个 \((i,j)\in\mathrm{I}^{2}\) ，存在整数 q 使得 \(1\leqslant q\leqslant m\) 且 \(a_{ij}^{(q)}>0\) 。（设 r 是 \((1+\mathrm{X})^{n-1}\) 除以 \(\psi\) 的欧几里得除法的余式；考察矩阵 \(r(A)\) 并注意它是严格正的。）

\begin{enumerate}
\def\labelenumi{\arabic{enumi})}
\setcounter{enumi}{11}
\tightlist
\item
  设 A 是 (I, I) 型的不可约且正的实矩阵。对每个非零的 \(x \in R^{I}\) ，我们置
\end{enumerate}

\[
r_{x} = \inf_{\substack{i\in \mathrm{I}\\ x_{i}\neq 0}}\frac{(Ax)_{i}}{x_{i}},
\]

其中 \((Ax)_i = \sum_{k\in I}a_{ik}x_k\) 是 \(Ax\) 的以 \(i\) 为指标的坐标。

\begin{enumerate}
\def\labelenumi{\alph{enumi})}
\item
  若 \(x \succcurlyeq 0\) 且 \(x \neq 0\) ，则 \(r_x\) 是最大的实数 \(\varrho\) ，使得 \(\varrho x \preccurlyeq Ax\) 。
\item
  假设 \(x \succcurlyeq 0\) ，并置 \(y = (1_{\mathrm{I}} + A)x\) 。则有 \(r_x \leqslant r_y\) 。
\item
  限制在 \((\mathbf{R}_{+}^{*})^{\mathrm{I}}\) 上时，函数 \(x \mapsto r_{x}\) 是连续的。
\end{enumerate}

\(d)\) 由此推出存在 \(z \in \mathbf{R}^{\mathrm{I}}\) 满足 \(z \succcurlyeq 0\) 且 \(r_z = \sup_{x \succcurlyeq 0} r_x\) （利用紧集 \(\mathrm{M} = \{x \in \mathbf{R}_{+}^{\mathrm{I}} \mid \sum_{i \in \mathrm{I}} x_i = 1\}\) 以及 \((1_{\mathrm{I}} + A)^{n-1}(\mathrm{M})\) ，再用习题 10。）

我们置 \(r = r_{z}\) 。\(R^{I}\) 中满足 \(z \succcurlyeq 0\) 且 \(r_{z} = r\) 的元素 z 称为极端元素。

\begin{enumerate}
\def\labelenumi{\alph{enumi})}
\setcounter{enumi}{4}
\item
  我们有 r \textgreater{} 0 。
\item
  若 z 是极端向量，则 Az = rz。（若 \(u = Az - rz \succcurlyeq 0\) 非零，则由习题 10，\((1_{\mathrm{I}} + A)^{n-1}u\) 将为 \(\succ 0\) ，且对 \(x = (1_{\mathrm{I}} + A)^{n-1}z\) 将有 \(Ax \succ rx\) 。）
\item
  每个极端元素都 \(\succ 0\) 。
\end{enumerate}

对每个向量 \(y\) （在 \(\mathbf{C}^{\mathrm{I}}\) 中），用 \(y^{+}\) 表示向量 \((|y_{i}|)_{i\in \mathrm{I}}\) ；于是有 \(y^{+}\in \mathbf{R}^{\mathrm{I}}\) 且 \(y^{+}\succcurlyeq 0\) 。

\begin{enumerate}
\def\labelenumi{\alph{enumi})}
\setcounter{enumi}{7}
\item
  设 \(\alpha \in \mathbf{C}\) 是把 \(A\) 视为复矩阵时的一个特征值，\(y\) 是 \(\mathbf{C}^{\mathrm{I}}\) 中相应的一个特征向量。我们有 \(|\alpha|y^{+} \preccurlyeq Ay^{+}\) 且 \(|\alpha| \leqslant r\) 。
\item
  推出：若 \(\alpha = r\) ，则 y 的所有坐标都非零，进而 A 的属于特征值 r 的特征空间（在 \(C^{I}\) 中）的维数为 1。
\end{enumerate}

用 \(\Delta(\mathrm{X})\) 表示 \(\mathrm{X} \cdot 1_{\mathrm{I}} - A\) 的行列式，用 \(B(\mathrm{X})\) 表示 \(\mathrm{X} \cdot 1_{\mathrm{I}} - A\) 的余子式矩阵的转置。我们记 \(B(\mathrm{X}) = (B_{ij}(\mathrm{X}))_{i,j \in \mathrm{I}}\) 。于是有 \(B(\mathrm{X})(\mathrm{X} \cdot 1_{\mathrm{I}} - A) = (\mathrm{X} \cdot 1_{\mathrm{I}} - A)B(\mathrm{X}) = \Delta(\mathrm{X}) \cdot 1_{\mathrm{E}}\) （A, III, p. 99, 命题 12）。

\begin{enumerate}
\def\labelenumi{\alph{enumi})}
\setcounter{enumi}{9}
\tightlist
\item
  对每个 \(i \in I\) ，\(B(r)\) 的以 i 为指标的列非零。（若不然，构造 A 的属于特征值 r 的一个特征向量，其第 i 个系数为零，这与 f) 矛盾。）
\end{enumerate}

\(k\) ) \(B(r)\) 的所有系数都有相同的符号。（首先利用关系式 \((r1_{\mathrm{I}} - A)B(r) = 0\) 证明：对所有 \(i\in \mathrm{I}\) ，以 \(i\) 为指标的 \(B(r)\) 的列的所有系数非零且同号，然后再对转置进行论证。）

\(l)\) \(\Delta'(r) > 0\) 且 \(B_{ij}(r) > 0\) 对所有 \((i,j) \in \mathrm{I}^2\) 成立（利用公式 \(\Delta'(X) = \sum_i B_{ii}(X)\) ，并注意 \(\Delta'\) 在区间 \([r, +\infty]\) 上不变号。）

\begin{enumerate}
\def\labelenumi{\alph{enumi})}
\setcounter{enumi}{12}
\tightlist
\item
  由此得出 r 是 A 的特征多项式的单根。
\end{enumerate}

\begin{enumerate}
\def\labelenumi{\arabic{enumi})}
\setcounter{enumi}{12}
\tightlist
\item
  设 A（相应地，C）是 (I, I) 型的实（相应地，复）方阵。假设 A 是正的且不可约，并且 \(C^{+} \preccurlyeq A\) ，这里与习题 12 一样，用 \(C^{+}\) 表示其系数为 C 的系数的绝对值的矩阵，而记号 \(C^{+} \preccurlyeq A\) 表示 \(A - C^{+}\) 的系数都是正的。
\end{enumerate}

我们沿用上一习题的记号。

\begin{enumerate}
\def\labelenumi{\alph{enumi})}
\item
  设 \(\gamma \in \mathbf{C}\) 是 \(C\) 的一个特征值，\(y\) 是 \(\mathbf{C}^{\mathrm{I}}\) 中相应的一个特征向量。我们有 \(|\gamma|y^{+} \preccurlyeq C^{+}y^{+} \preccurlyeq Ay^{+}\) 且 \(|\gamma| \leqslant r_{y^{+}} \leqslant r\) 。
\item
  我们进一步假设 \(|\gamma| = r\) 。则 \(y^{+}\) 是 A 的一个极端向量（参见习题 12 的问题 f)），且有关系式 \(Ay^{+} = C^{+}y^{+} = ry^{+}\) 与 \(A = C^{+}\) 。
\end{enumerate}

设 \(u\)（相应地，\(u_j\) ，\(j \in I\) ）是使 \(\gamma = |\gamma|u\) （相应地，\(y_j = |y_j|u_j\) ，\(j \in I\) ）成立的唯一复数。用 \(D\) 表示 (I, I) 型的复方阵，其对角系数为 \((u_j)_{j \in I}\) ，其余系数为零，于是 \(y = Dy^+\)。设 \(F = u^{-1}D^{-1}CD\) 。

\begin{enumerate}
\def\labelenumi{\alph{enumi})}
\setcounter{enumi}{2}
\item
  有 \(Fy^{+} = ry^{+}\) 且 \(F^{+}y^{+} = Fy^{+}\) （注意 \(F^{+} = C^{+}\) 。）
\item
  有 \(F = F^{+}\) 且 \(C = uDAD^{-1}\) 。
\end{enumerate}

\begin{enumerate}
\def\labelenumi{\arabic{enumi})}
\setcounter{enumi}{13}
\tightlist
\item
  设 A 是 (I, I) 型的正且不可约的实方阵。我们沿用习题 12 的记号。
\end{enumerate}

设 \((\lambda_{h})_{h\in\mathrm{H}}\) 是 A 的模为 r 的复特征值按重数重复而得的族；置 \(v_{h}=\lambda_{h}/r\) （对所有 \(h\in H\) ）。

设 \(h \in \mathrm{H}\) ；把习题 13 应用于 \(C = A\) 与 \(\gamma = \lambda_h\) ，存在矩阵 \(D_{(h)}\) 使得 \(A = v_h D_{(h)} AD_{(h)}^{-1}\) 且 \(D_{(h)}^+ = 1_{\mathrm{I}}\) 。

\begin{enumerate}
\def\labelenumi{\alph{enumi})}
\item
  设 \(z \succ 0\) 是 A 的一个极端向量。对 \(h \in H\) ，向量 \(y_{(h)} = D_{(h)}z\) 是 A 的属于特征值 \(\lambda_{h}\) 的特征向量。
\item
  证明：对 \(h \in H\) ，A 的特征值 \(\lambda_{h}\) 是单重的。由此推出矩阵 \(D_{(h)}\) 与向量 \(y_{(h)}\) 在相差一个纯量因子的意义下唯一确定。
\item
  证明：对 \(H^{2}\) 中所有 h 与 k，有 \(A = u_{h}u_{k}^{-1}D_{(h)}D_{(k)}^{-1}AD_{(k)}D_{(h)}^{-1}\) ，从而向量 \(D_{(h)}D_{(k)}^{-1}z\) 是 A 的属于特征值 \(u_{h}u_{k}^{-1}r\) 的特征向量。
\item
  由此推出 \(\{u_h\}_{h\in \mathrm{H}}\) 是 \(\mathbf{C}^*\) 的一个子群，且存在非零复数的一个族 \((\mu_h)_{h\in \mathrm{H}}\) ，使得 \(\{\mu_h^{-1}D_{(h)}\}_{h\in \mathrm{H}}\) 是 (I,I) 型复系数可逆矩阵所成的群的一个子群，并且这两个子群都是阶为 Card(H) 的循环群。（选取一个指标 \(i\in \mathrm{I}\) ，并把 \(\mathrm{D}_{(h)}\) 规范化，使其行指标与列指标均为 \(i\) 的系数等于 1。）
\end{enumerate}

在本习题的其余部分，我们假设 H = Z/dZ，且存在 \(\omega \in C^{*}\) 使得 \(\omega^{d} = 1\) 且 \(u_{h} = \omega^{h}\) （\(h \in H\) ）。置 \(D = D_{(1)}\) ，则 \(D^{d} = 1_{I}\) 且 \(D_{(h)} = D^{h}\) 对所有 \(h \in H\) 成立。

\begin{enumerate}
\def\labelenumi{\alph{enumi})}
\setcounter{enumi}{4}
\tightlist
\item
  由此推出 \(A = \omega DAD^{-1}\) 。
\end{enumerate}

我们考虑 I 的划分 \((\mathrm{I}_{h})_{h\in\mathrm{H}}\) ，使得对 \(i\in I\) ，D 的以 \((i,i)\) 为指标的系数等于 \(\omega^{h}\) 当且仅当 \(i\in I_{h}\) 。对每个 h，置 \(n_{h}=\mathrm{Card}(I_{h})\) 。矩阵 D 写成“分块”形式

\[
D = \left( \begin{array}{c c c c} 1 _ {n _ {0}} & & & \\ & \omega 1 _ {n _ {1}} & & \\ & & \ddots & \\ & & & \omega^ {d - 1} 1 _ {n _ {d - 1}} \end{array} \right)
\]

我们考虑 \(A\) 的相应的分解：

\[
A = \left( \begin{array}{c c c c} A _ {0, 0} & A _ {0, 1} & \dots & A _ {0, d - 1} \\ A _ {1, 0} & A _ {1, 1} & \dots & A _ {1, d - 1} \\ \vdots & \vdots & \ddots & \vdots \\ A _ {d - 1, 0} & A _ {d - 1, 1} & \dots & A _ {d - 1, d - 1} \end{array} \right),
\]

使得对 \(h, k \in \mathrm{H}\) ，\(A_{h,k}\) 是 \((\mathrm{I}_h, \mathrm{I}_k)\) 型的矩阵。

\(f)\) 对所有 \((h,k)\in\mathrm{H}\) ，有 \(\omega A_{h,k}=\omega^{k-h}A_{h,k}\) 。由此推出 \(\mathrm{A}_{h,k}=0\) （若 \(k-h\neq1\) ）。

\begin{enumerate}
\def\labelenumi{\arabic{enumi})}
\setcounter{enumi}{14}
\tightlist
\item
  我们沿用习题 12 的假设与记号。设 \(\Delta_1(\mathrm{X})\) 是 \(B(\mathrm{X})\) 的诸系数的最大公因子；假设多项式 \(\Delta_1(\mathrm{X})\) 是首一的。
\end{enumerate}

\begin{enumerate}
\def\labelenumi{\alph{enumi})}
\item
  多项式 \(\Delta_1(X)\) 整除 \(\Delta(X)\) ，且 \(\Delta_1(r) \neq 0\) 。
\item
  由此推出 \(\Delta_{1}(X)\) 的实根都严格小于 r。从而得出 \(\Delta_{1}(r)>0\) 。
\item
  置 \(C(\mathrm{X}) = \frac{1}{\Delta_1(\mathrm{X})} B(\mathrm{X})\) 。这是一个系数属于 \(\mathbf{R}[\mathrm{X}]\) 的矩阵。证明 \(C(r) \succ 0\) 。
\end{enumerate}

\begin{enumerate}
\def\labelenumi{\arabic{enumi})}
\setcounter{enumi}{15}
\tightlist
\item
  设 A 是 (I, I) 型的正且不可约的实方阵。对所有 \(j \in I\) 置 \(s_{j} = \sum_{k \in I} a_{jk}\) ，
\end{enumerate}

\[
s = \inf _ {j \in \mathrm{I}} s _ {j}, \quad \text {and} \quad \mathrm{S} = \sup _ {j \in \mathrm{I}} s _ {j}.
\]

我们也采用习题 12 的记号 \(r\) 与 \(B(\mathbf{X})\) 。

\begin{enumerate}
\def\labelenumi{\alph{enumi})}
\item
  有 \(\sum_{j\in \mathrm{I}}(r - s_j)B_{kj}(r) = 0\)
\item
  由此推出 \(s \leqslant r \leqslant S\) ，且这两个不等式中任何一个取等号都蕴含 s = S。
\end{enumerate}

\begin{enumerate}
\def\labelenumi{\arabic{enumi})}
\setcounter{enumi}{16}
\item
  设 \(A\) 是 (I, I) 型的正且不可约的实方阵。设 \(y \in \mathbf{R}^{\mathrm{I}}\) 是 \(A\) 的属于特征值 \(\alpha\) 的特征向量，且满足 \(y \succcurlyeq 0\) 。则 \(\alpha = r\) （沿用习题 12 的记号）且 \(y \succ 0\) 。（考虑 \(u\) （\(^{t}A\) 的一个极端向量），并计算 \(\langle y|^{t}Au\rangle = \langle Ay|u\rangle\) 。）
\item
  设 A 是 (I, I) 型的正且不可约的实方阵。设 \(\varrho\) 是 A 的谱半径。对每个 \(x \in R^{I}\) （满足 \(y \succcurlyeq 0\) 且 \(x \neq 0\) ），我们置
\end{enumerate}

\[
\varrho^ {x} = \sup _ {i \in \mathrm{I}} \frac {(A x) _ {i}}{x _ {i}},
\]

（约定 \(\varrho^x = +\infty\) ，如果存在指标 \(i\) 使 \(x_i = 0\) 且 \((Ax)_i \neq 0\) ）。设

\[
\varrho^{\prime} = \inf_{\substack{x\succ 0\\ x\neq 0}}\varrho^{x}.
\]

沿用与习题 12 类似的途径，证明 \(\varrho'\) 是 \(A\) 的一个特征值且 \(\varrho' = \varrho\) 。

\begin{enumerate}
\def\labelenumi{\arabic{enumi})}
\setcounter{enumi}{18}
\tightlist
\item
  设 A 是 (I, I) 型的正且不可约的实方阵；\(\varrho\) 是它的谱半径。设 \(z \in R^{I}\) 满足 \(z \succcurlyeq 0\) 且 \(z \neq 0\) 。
\end{enumerate}

若 \(\varrho z \preccurlyeq Az\) （相应地，\(\varrho z \succcurlyeq Az\) ），则 \(\varrho z = Az\) 且 \(z \succ 0\) 。

\begin{enumerate}
\def\labelenumi{\arabic{enumi})}
\setcounter{enumi}{19}
\tightlist
\item
  设 A 是 (I, I) 型的正但未必不可约的实方阵；\(\varrho\) 是 A 的谱半径。我们保留前面诸习题的记号 \(B(\mathrm{X})\) 与 \(C(\mathrm{X})\) 。
\end{enumerate}

\begin{enumerate}
\def\labelenumi{\alph{enumi})}
\item
  实数 \(\varrho\) 是 \(A\) 的特征值，且存在至少一个特征向量 \(\succcurlyeq 0\) 相伴于 \(\varrho\) 。
\item
  对每个实数 \(\lambda \geqslant r\) ，矩阵 \(B(\lambda)\) 与 \(B'(\lambda)\) 都是正的。（对第二个矩阵，可对 I 的基数用归纳法。）
\item
  矩阵 \(C(\lambda)\) 对所有 \(\lambda \geqslant r\) 都是正的。
\item
  若 A 是不可约的，则对每个实数 \(\lambda > r\) ，矩阵 \(B(\lambda)\) 、\(C(\lambda)\) 与 \((\lambda1_{\mathrm{I}} - A)^{-1}\) 都是严格正的。
\item
  设 J 是 I 的异于 I 的子集。设 \(\varrho_{\mathrm{J}}\) 是子矩阵 \((a_{i,j})_{i,j\in \mathrm{J}}\) 的谱半径。我们有 \(\varrho_{\mathrm{J}}\leqslant \varrho\) ，且当 \(A\) 不可约时该不等式是严格的。反之，若 \(A\) 是可约的，则存在异于 I 的子集 \(\mathrm{J}\subset \mathrm{I}\) 使得 \(\varrho_{\mathrm{J}} = \varrho\) 。
\item
  设 J 是 I 的异于 I 的子集，且子矩阵 \((a_{i,j})_{i,j\in\mathrm{J}}\) 的谱半径等于 \(\varrho\) 。对每个满足 \(J\subset K\subset I\) 的子集 K，矩阵 \((a_{i,j})_{i,j\in\mathrm{K}}\) 的谱半径等于 \(\varrho\) 。由此推出 \(A\succ0\) 当且仅当 \(B(r)\) 的所有对角系数都严格为正。
\end{enumerate}

\(g)\) 对每个实数 \(\lambda >\varrho\) 以及 I 的每个满足 \(\mathrm{J}\neq \mathrm{I}\) 的子集 J，矩阵 \((\lambda \delta_{i,j} - a_{i,j})_{i,j\in \mathrm{J}}\) 的行列式严格为正。

\begin{enumerate}
\def\labelenumi{\alph{enumi})}
\setcounter{enumi}{7}
\tightlist
\item
  设 \(\lambda\) 是实数；假设 \((\lambda \delta_{i,j} - a_{i,j})_{i,j \in \mathbf{J}}\) 的行列式对每个子集 \(\mathbf{J} \subset \mathbf{I}\) （异于 \(\mathbf{I}\) ）都严格为正。则有 \(\lambda > \varrho\) 。
\end{enumerate}

在下面的问题中，我们假设 \(\mathrm{I}\) 是区间 \([1,n]\) （在 \(\mathbf{N}\) 中）。i) 设 \((g_{i,j})_{1\leqslant i,j\leqslant n}\) 是一个实矩阵，使得 \(g_{i,j}\leqslant0\) 对所有 \(i\neq j\) 成立，且矩阵 \((g_{i,j})_{1\leqslant i,j\leqslant m}\) 的行列式对每个整数 \(m\) （满足 \(1\leqslant m\leqslant n\) ）都严格为正。则对 I 的每个子集 J 有 \(\det(g_{i,j})_{i,j\in\mathrm{J}}>0\) 。

\begin{enumerate}
\def\labelenumi{\alph{enumi})}
\setcounter{enumi}{9}
\item
  设 \(\lambda\) 是实数；假设对每个整数 \(m\) （满足 \(1 \leqslant m \leqslant n\) ），行列式 \(\det(\lambda \delta_{i,j} - a_{i,j})_{1 \leqslant i,j \leqslant m}\) 严格为正。则 \(\lambda > \varrho\) 。
\item
  设 \(C = (c_{i,j})_{1 \leqslant i,j \leqslant n}\) 是一个实矩阵，使得 \(c_{i,j} \geqslant 0\) 对所有 \(i \neq j\) 成立。若 \((-1)^{k} \det(c_{i,j})_{1 \leqslant i,j \leqslant m}\) 对每个满足 \(1 \leqslant m \leqslant n\) 的整数 m 都严格为正，则 C（视为复矩阵）的特征值的实部都是负的。
\end{enumerate}

\begin{enumerate}
\def\labelenumi{\arabic{enumi})}
\setcounter{enumi}{20}
\tightlist
\item
  设 A 是 (I, I) 型的正的实方阵；\(\varrho\) 是它的谱半径。
\end{enumerate}

\begin{enumerate}
\def\labelenumi{\alph{enumi})}
\tightlist
\item
  证明：存在整数 p 与 q 满足 \(0 \leqslant q \leqslant p\) ，以及 I 的一个分成非空部分的划分 \((\mathrm{I}_{s})_{1 \leqslant s \leqslant p}\) ，使得相应的分块分解 \(A = (A_{s,t})_{1 \leqslant s, t \leqslant p}\) 具有下列性质：
\end{enumerate}

- 矩阵 \(A_{s,s}\) 都是不可约的；

- 我们有 \(A_{s,t} = 0\) （若 \(t > s\) ）；

- 我们有 \(A_{s,t} = 0\) （若 \(1 \leqslant t < s \leqslant q\) ）；

- 对每个整数 \(s\) （满足 \(q \leqslant s \leqslant p\) ），存在 \(t\) 使得 \(A_{s,t} \neq 0\) 。

\begin{enumerate}
\def\labelenumi{\alph{enumi})}
\setcounter{enumi}{1}
\item
  确定这种划分的所有可能情形。（特别地，证明整数 q 与 p 都是唯一确定的，集组 \(\{I_{1},\ldots,I_{g}\}\) 与集组 \(\{I_{q+1},\ldots,I_{p}\}\) 亦然。）
\item
  存在 A 的一个特征向量 \(z \succ 0\) 属于特征值 \(\varrho\) ，当且仅当对每个整数 \(s \in [1, p]\) ，\(A_{s,s}\) 的谱半径等于 \(\varrho\) （若 \(s \leqslant q\) ），否则严格小于 \(\varrho\) 。
\item
  为使 \(A\) 与 \({}^t A\) 二者都承认一个严格正的极端向量，必须且只需存在如上的一个划分，使得 \(A_{s,t} = 0\) （对 \(1 \leqslant t < s\) 且 \(q < s \leqslant p\) ），且矩阵 \(A_{s,s}\) （\(s \leqslant q\) ）的谱半径等于 \(\varrho\) 。
\item
  若 A 与 \(^{t}A\) 都承认一个严格正的极端向量，且 \(\varrho\) 的谱重数等于 1，则 A 是不可约的。
\end{enumerate}

\begin{enumerate}
\def\labelenumi{\arabic{enumi})}
\setcounter{enumi}{21}
\tightlist
\item
  设 A 是 (I, I) 型的正且不可约的实方阵；\(\varrho\) 是它的谱半径。设 \(d \geqslant 1\) 是由习题 14 确定的整数。当 d = 1 时，我们称 A 是本原的。设 \(\Delta(\mathrm{X})\) 是多项式 \(\det(\mathrm{X} \cdot 1_{\mathrm{I}} - A)\) 。
\end{enumerate}

\begin{enumerate}
\def\labelenumi{\alph{enumi})}
\item
  我们写 \(\Delta(\mathrm{X}) = \mathrm{X}^{n} + a_{1}\mathrm{X}^{n_{1}} + \cdots + a_{t}\mathrm{X}^{n_{t}}\) ，其中 \(a_{1}, \ldots, a_{t}\) 是非零实数，\(n_{1}, \ldots, n_{t}\) 是满足 \(n > n_{1} > \cdots > n_{t}\) 的自然数。证明 d 是 \(n - n_{1}, n_{1} - n_{2}, \ldots, n_{t-1} - n_{t}\) 的最大公因子。
\item
  存在非零整数 m 与多项式 g，使得 \(\Delta(\mathrm{X}) = g(\mathrm{X}^{d})\mathrm{X}^{m}\) 。
\item
  若存在整数 p \textgreater{} 0 使得 \(A^{p}\) 严格正，则 A 是本原的。
\item
  假设 \(A\) 是本原的。引入习题 15 的矩阵 \(C(\mathbf{X})\) ，并按该习题的记号置 \(\mathsf{X}(\mathbf{X}) = \Delta (\mathbf{X}) / \Delta_1(\mathbf{X})\) 。序列 \((A^k /\varrho^k)_{k\in \mathbf{N}}\) 的极限等于 \(C(\varrho) / \mathsf{X}'(\varrho)\) 。（可利用 \(A\) 的若尔当块分解。）利用习题 15 的结果，推出存在整数 \(p > 0\) 使得 \(A^p\) 严格正。
\item
  假设存在整数 q \textgreater{} 0 使得 \(A^{q}\) 是可约的。设 p 是 q 与 d 的最大公因子。证明：存在 I 的一个分成非空部分的划分 \((\mathrm{I}_{s})_{1 \leqslant s \leqslant p}\) ，使 \(A^{q}\) 关于此划分是分块对角的，其对角块不可约且谱半径为 \(\varrho^{q}\) 。
\item
  若 A 是本原的，则 \(A^{p}\) 对每个 p \textgreater{} 0 都是不可约的。
\end{enumerate}

\(g)\) 存在 I 的一个分成非空部分的划分 \((\mathrm{I}_s)_{1 \leqslant s \leqslant d}\) ，使 \(A^d\) 关于此划分是分块对角的，其对角块是本原的且谱半径等于 \(\varrho^d\) 。

\begin{enumerate}
\def\labelenumi{\arabic{enumi})}
\setcounter{enumi}{22}
\tightlist
\item
  设 P 是 (I, I) 型的实方阵；若 P 是非负的，且向量 \((1, \ldots, 1) \in \mathbf{R}^{\mathrm{I}}\) 是 P 的属于特征值 1 的特征向量，则称 P 是随机的。
\end{enumerate}

设 A 是 (I, I) 型的实方阵。证明：A 是正的且承认一个严格正的极端向量，当且仅当存在一个随机矩阵 P、一个实数 r \textgreater{} 0 以及一个对角系数严格为正的对角矩阵 Z，使得 \(A = rZPZ^{-1}\) 。在这种情形，\(A\) 的谱半径等于 \(r\) ，而 \(A\) 的一个极端向量是 \(Z(1, \ldots, 1)\) 。

\begin{enumerate}
\def\labelenumi{\arabic{enumi})}
\setcounter{enumi}{23}
\tightlist
\item
  \begin{enumerate}
  \def\labelenumii{\alph{enumii})}
  \tightlist
  \item
    设 P 是 n 阶的实随机方阵。实数 1 是 P 的极小多项式的单根。（可利用习题 20。）此外，P 的极小多项式的每个模为 1 的根都是单根。
  \end{enumerate}
\end{enumerate}

\begin{enumerate}
\def\labelenumi{\alph{enumi})}
\setcounter{enumi}{1}
\tightlist
\item
  设 A 是一个正实矩阵，它承认一个极端向量 \(z \succ 0\) ，谱半径为 \(\varrho\) 。A 的极小多项式的所有模为 \(\varrho\) 的根都是单根。
\end{enumerate}

在习题 25 至 30 中，我们考虑一个实弗雷歇空间 E 和 E 中的一个闭凸锥 C，并假设 C 的内部非空。由 C 诱导的 E 上的序关系记作 \(\preccurlyeq_{\mathrm{C}}\) 。于是，对 E 中的 \(x\) 、\(y\) ，\(x \succcurlyeq_{\mathrm{C}} y\) 当且仅当 \(x - y \in \mathrm{C}\) 。我们将用 \(x \succ_{\mathrm{C}} y\) 表示关系 \(x - y \in \mathring{\mathrm{C}}\) 。

我们考虑满足下列条件的连续映射 u: E → E：

\begin{enumerate}
\def\labelenumi{(\roman{enumi})}
\item
  映射 u 是正齐次的：\(u(\lambda x) = \lambda u(x)\) 对每个实数 \(\lambda \geqslant 0\) 成立；
\item
  映射 u 关于 C 是递增的，即 \(x \succsim_{C} y\) 蕴含 \(u(x) \succsim_{C} u(y)\) ；
\item
  对每个 \(x \in \mathrm{C} - \{0\}\) ，有 \(u(x) \in \mathring{\mathrm{C}}\) 。
\end{enumerate}

这些习题取自 P.-L. 利翁斯 2020/2021 学年在法兰西公学院的课程，其目标是在某些条件下证明：存在元素 \(x \in \mathrm{C} - \{0\}\) 与 \(\lambda > 0\) 使得 \(x = \lambda u(x)\) ；特别地，这给出了克赖因–鲁特曼定理（III, p. 94, 定理 4 与 III, p. 97, 命题 8）的另一个证明。

对 \(f \in \mathring{\mathrm{C}}\) 与 \(\lambda \in \mathbf{R}_{+}\) ，我们把方程 \(y = \lambda u(y) + f\) 的上解（相应地，下解）称为一个元素 \(y \in \mathrm{C}\) ，它满足 \(y \succcurlyeq_{\mathrm{C}} \lambda u(y) + f\) （相应地，满足 \(y \preccurlyeq_{\mathrm{C}} \lambda u(y) + f\) ）。我们用 \(U_{f}\) 表示由 \(\lambda \in \mathbf{R}_{+}\) （使方程 \(y = \lambda u(y) + f\) 的上解集合非空）所成的集。

\begin{enumerate}
\def\labelenumi{\arabic{enumi})}
\setcounter{enumi}{24}
\tightlist
\item
  \begin{enumerate}
  \def\labelenumii{\alph{enumii})}
  \tightlist
  \item
    证明：对每个 \(x \in C\) 与每个 \(y \in \mathring{C}\) ，存在实数 \(\lambda \geqslant 0\) 使得 \(x \preccurlyeq_{C} \lambda y\) 。
  \end{enumerate}
\end{enumerate}

\begin{enumerate}
\def\labelenumi{\alph{enumi})}
\setcounter{enumi}{1}
\tightlist
\item
  集合 \(U_{f}\) 是包含 0 的一个区间。
\end{enumerate}

我们对映射 u 引入下列条件。

\begin{enumerate}
\def\labelenumi{(\roman{enumi})}
\setcounter{enumi}{3}
\item
  对每个 \(\lambda \in \mathbf{R}_{+}\) 、每个 \(f \in \mathring{\mathrm{C}}\) ，以及每个序列 \((x_{n})\) （由 C 中元素组成，满足 \(x_{n+1} = \lambda u(x_{n}) + f\) （对所有 \(n\) ），且递增并上有界或递减），它都是收敛的；
\item
  E 的每个有界子集在 u 下的像是相对紧的。
\item
  对每个映射 \(\lambda \mapsto \varepsilon_{\lambda}\) （从 \(U_f\) 到 C 中），若它满足 \(\lim_{\lambda \to \lambda_m} \varepsilon_{\lambda} = 0\) ，则集合
\end{enumerate}

\[
\{z \in \mathrm{C} \mid \exists \lambda \in \mathrm{U}_{f}\text{ 使得 } z = \lambda u(z) + \varepsilon_{\lambda} \}
\]

在 E 中是相对紧的。

\begin{enumerate}
\def\labelenumi{\arabic{enumi})}
\setcounter{enumi}{25}
\tightlist
\item
  我们沿用前面的记号与假设，并另外假设条件 (iv) 成立。
\end{enumerate}

\begin{enumerate}
\def\labelenumi{\alph{enumi})}
\item
  设 \(\lambda \in \mathrm{U}_f\) 。我们定义一个序列 \((x_n)_{n \in \mathbf{N}}\) ：\(x_0 = 0\) 且 \(x_{n+1} = \lambda u(x_n) + f\) （对所有 \(n \in \mathbf{N}\) ）。它递增且上有界；我们有 \(x_n \in \mathring{\mathrm{C}}\) （对 \(n \geqslant 1\) ）。
\item
  由条件 (iv)，序列 \((x_{n})\) 收敛；用 \(x(\lambda)\) 表示它的极限。证明 \(x(\lambda) \in \mathring{\mathrm{C}}\) 与 \(x(\lambda) = \lambda u(x(\lambda)) + f\) 。
\item
  设 \(\lambda \in \mathrm{U}_f\) ，并设 \(y \in \mathrm{C}\) 满足 \(y \succcurlyeq_{\mathrm{C}} \lambda u(y) + f\) ；由 \(x_0 = y\) 、\(x_{n+1} = \lambda u(x_n) + f\) 所定义的序列是递减的，取值于 C 中，其极限 \(x\) 满足方程 \(x = \lambda u(x) + f\) 。
\end{enumerate}

\(d)\) \(\mathrm{U}_{f}\) 在 \(\mathbf{R}_{+}\) 中的上确界不属于 \(\mathrm{U}_{f}\) 。

\begin{enumerate}
\def\labelenumi{\alph{enumi})}
\setcounter{enumi}{4}
\item
  \(U_{g} = U_{f}\) 对所有 \(g \in \mathring{C}\) 成立（对每个 \(\lambda \in U_{f}\) ，设 \(x \in C\) 满足 \(x \succsim_{C} \lambda u(x) + f\) ；则存在 M \textgreater{} 0 使得 y = Mx 满足 \(y \succsim_{C} \lambda u(y) + g\) 。） f) 设 \(\lambda \leqslant \mu\) 是实数。设 \(x\) 是方程 \(y = \lambda u(y) + f\) 的一个下解，设 \(x'\) 是方程 \(y = \mu u(y) + f\) 的一个上解。则 \(x \preccurlyeq_{C} x'\) 。（证明由 \(\varepsilon \in [0,1]\) （使 \(\varepsilon x \preccurlyeq_{C} x'\) ）组成的集与 {[}0,1{]} 重合。）
\item
  若 C 是尖锥，则对每个 \(\lambda \in \mathrm{U}_f\) ，存在唯一的元素 \(x \in \mathring{\mathrm{C}}\) 使得 \(x = \lambda u(x) + f\) 。
\end{enumerate}

\(h\) ) 设 \(\mu\in\mathbf{R}_{+}\) 满足 \(f\preccurlyeq_{\mathrm{C}}\mu u(f)\) 。则 \(\mu\notin\mathrm{U}_{f}\) 。（序列 \((x_{n})_{n\in\mathbf{N}}\) （由 \(x_{0}=0\) 与 \(x_{n+1}=\mu u(x_{n})+f\) 定义）满足 \(x_{n+1}-x_{n}\succcurlyeq_{\mathrm{C}}f\) （对所有 \(n\in\mathbf{N}\) ）。）特别地，区间 \(\mathrm{U}_{f}\) 有界。此后我们用 \(\lambda_{\mathrm{m}}\) 表示 \(\mathrm{U}_{f}\) 的上确界。

\begin{enumerate}
\def\labelenumi{\roman{enumi})}
\tightlist
\item
  设 \(z \in \mathrm{C} - \{0\}\) 与 \(\mu \in \mathbf{R}_+\) 满足 \(\mu u(z) \succsim_{\mathrm{C}} z\) 。则 \(\lambda_{\mathrm{m}} \leqslant \mu\) ，从而
\end{enumerate}

\[
\lambda_ {\mathrm{m}} \leqslant \inf _ {z \in \mathrm{C} - \{0 \}} \inf \{\mu \mid z \leqslant \mu u (z) \}.
\]

\begin{enumerate}
\def\labelenumi{\alph{enumi})}
\setcounter{enumi}{9}
\tightlist
\item
  映射 \(\lambda \mapsto x(\lambda)\) （其中 \(x(\lambda)\) 在问题 \(a\) 中定义）是递增的。
\end{enumerate}

\(k)\) 映射 \(\lambda \mapsto x(\lambda)\) 当 \(\lambda\) 趋于 \(\lambda_{\mathrm{m}}\) 时没有极限。更确切地说，若 \((\lambda_n)_n\) 是 \(\mathrm{U}_f\) 中趋于 \(\lambda_{\mathrm{m}}\) 的一个序列，则序列 \((x(\lambda_n))_n\) 不收敛。

\begin{enumerate}
\def\labelenumi{\arabic{enumi})}
\setcounter{enumi}{26}
\tightlist
\item
  我们假设 (i) 至 (v) 各条假设成立。
\end{enumerate}

\begin{enumerate}
\def\labelenumi{\alph{enumi})}
\item
  证明：\(x(\lambda)\) （对 \(\lambda \in \mathrm{U}_f\) ）所成的集不是有界的。
\item
  设 d 是 E 上定义其拓扑的一个距离。对 \(\lambda \in R_{+}\) ，定义 \(\hat{x}(\lambda) = d(0, x(\lambda))^{-1}x(\lambda)\) 。证明映射 \(\hat{x}\) 有界，且 \(\hat{x}\) 的值集有一个聚点 \(x_{m}\) 满足
\end{enumerate}

\begin{enumerate}
\def\labelenumi{(\roman{enumi})}
\item
  \(d(0, x_{\mathrm{m}}) = 1,\)
\item
  \(x_{\mathrm{m}} = \lambda_{\mathrm{m}} u(x_{\mathrm{m}}),\)
\item
  \(x_{m} \in \overset{\circ}{C}.\)
\end{enumerate}

\begin{enumerate}
\def\labelenumi{\arabic{enumi})}
\setcounter{enumi}{27}
\tightlist
\item
  在本习题中我们假设条件 (i) 至 (iv) 成立。 a) 我们固定 \(x_0 \in \mathring{\mathrm{C}}\) 。对 \(\lambda \in \mathrm{U}_f\) ，令 \(\mathrm{A}(\lambda) = \inf \{\mathrm{M} \in \mathbf{R}_+ \mid x(\lambda) \preccurlyeq_{\mathrm{C}} \mathrm{M}x_0\}\) 。设 \(\mathrm{C}_0 \in \mathbf{R}_+\) 满足 \(f \preccurlyeq_{\mathrm{C}} \mathrm{C}_0 u(x_0)\) 。对每个 \(\lambda\) （\(\mathrm{U}_f\) 中的），区间 \([\lambda, \lambda + \mathrm{C}_0 / \mathrm{A}(\lambda)[\) 包含于 \(\mathrm{U}_f\) （参见习题 26 的问题 \(d\) )）。
\end{enumerate}

\begin{enumerate}
\def\labelenumi{\alph{enumi})}
\setcounter{enumi}{1}
\tightlist
\item
  由此推出：函数 \(\lambda \mapsto \mathrm{A}(\lambda)\) 趋于 \(+\infty\) （当 \(\lambda \in \mathrm{U}_f\) 趋于 \(\lambda_{\mathrm{m}}\) 时）。
\end{enumerate}

对每个 \(\lambda\) （\(\mathrm{U}_f\) 中的），定义 \(\hat{x} (\lambda) = x(\lambda) / \mathrm{A}(\lambda)\) 。

我们现在假设条件 (vi) 成立。

\begin{enumerate}
\def\labelenumi{\alph{enumi})}
\setcounter{enumi}{2}
\tightlist
\item
  证明映射 \(\hat{x}\) 的值集有一个聚点 \(x_{\mathrm{m}}\) 满足
\end{enumerate}

\[
0 \preccurlyeq_ {\mathrm{C}} x _ {\mathrm{m}} \preccurlyeq_ {\mathrm{C}} x _ {0} \quad \text {and} \quad x _ {\mathrm{m}} = \lambda_ {\mathrm{m}} u (x _ {\mathrm{m}}).
\]

\begin{enumerate}
\def\labelenumi{\alph{enumi})}
\setcounter{enumi}{3}
\item
  证明 \(x_{m} \neq 0\) 。（若 \(x_{m} = 0\) ，设 U 是 E 中 0 的一个邻域，使得 \(x_{0} + U \subset C\) ；则对每个 \(\mu \geqslant 0\) ，将存在一个序列 \((\lambda_{n})\) （\(U_{f}\) 中的），收敛于 \(\lambda_{m}\) ，使 \(x_{0} - \mu \hat{x}(\lambda_{n}) \succsim_{C} 0\) ，这与 \(A(\lambda)\) 的定义矛盾。）
\item
  设 \(y \in \mathrm{C} - \{0\}\) 满足 \(y = \lambda_{\mathrm{m}} u(y)\) 。置 \(\theta = \sup \{\gamma \mid \gamma y \preccurlyeq_{\mathrm{C}} x_{\mathrm{m}}\}\) 。我们有 \(\theta \in \mathbf{R}_{+}^{*}\) 且 \(x_{\mathrm{m}} = \theta y\) 。（我们有 \(\theta y \preccurlyeq_{\mathrm{C}} x_{\mathrm{m}}\) 与 \(\theta y = \lambda_{\mathrm{m}} u(\theta y)\) ；若 \(x_{\mathrm{m}} \neq \theta y\) ，则将有 \(x_{\mathrm{m}} - \theta y \succ_{\mathrm{C}} 0\) ，且存在 \(\varepsilon > 0\) 使 \(x_{\mathrm{m}} - \theta y \succcurlyeq_{\mathrm{C}} \varepsilon y\) ，这与 \(\theta\) 的定义矛盾。）
\item
  设 \(y \in C - \{0\}\) 与 \(\mu \in R\) 满足 \(y = \mu u(y)\) 。则 \(\mu = \lambda_{m}\) 。（注意 \(0 < \mu\) 且 \(\mu \geqslant \lambda_{m}\) ；若有 \(\mu > \lambda_{m}\) ，令 \(\theta = \sup\{\gamma \mid \gamma x_{m} \preccurlyeq_{C} y\}\) ；则将有 \(\theta \in R_{+}^{*}\) ，进而 \(\theta x_{m} \preccurlyeq_{C} y\) ，于是 \(\theta x_{m} = \lambda_{m} u(\theta x_{m}) \prec_{C} \mu u(\theta x_{m}) \preccurlyeq_{C} \mu u(y) = y\) ，这与 \(\theta\) 的定义矛盾。）
\end{enumerate}

\begin{enumerate}
\def\labelenumi{\arabic{enumi})}
\setcounter{enumi}{28}
\tightlist
\item
  我们沿用前一习题的记号；特别地，假设 C 的内部非空。我们假设 u 是线性映射，且条件 (i)、(ii)、(iii) 成立。
\end{enumerate}

\begin{enumerate}
\def\labelenumi{\alph{enumi})}
\tightlist
\item
  若 u 是 E 的紧自同态，则条件 (iv) 与 (v) 成立。
\end{enumerate}

我们假设条件 (iv) 与 (vi) 成立。设 \((\lambda_{\mathrm{m}}, x_{\mathrm{m}})\) 是由前一习题给出的。

\begin{enumerate}
\def\labelenumi{\alph{enumi})}
\setcounter{enumi}{1}
\tightlist
\item
  对每个 \(\lambda \in \mathbf{R}_{+}\) （满足 \(\lambda < \lambda_{\mathrm{m}}\) ），自同态 \(1_{\mathrm{E}} - \lambda u\) 是可逆的。（首先证明 \(1_{\mathrm{E}} - \lambda u\) 通过过渡到子空间在 \(\mathring{\mathrm{C}}\) 上诱导一个双射，参见前一习题的问题 \(a\) ) 与 \(f\) )；然后推出
\end{enumerate}

\(1_{E} - \lambda u\) 是满射，然后通过把 \(x_{1} - x_{2}\) 写成 \(1_{E} - \lambda u\) 的核中元素（其中 \(x_{i} \in \mathring{C}\) ）并验证存在 \(z \in E\) 使 \(x_{1} + z - \lambda u(x_{1} + z) \in \mathring{C}\) ，推出它是单射。

\begin{enumerate}
\def\labelenumi{\alph{enumi})}
\setcounter{enumi}{2}
\tightlist
\item
  设 E 中的 y 满足 \(y = \lambda_{\mathrm{m}} u(y)\) 。则 \(y \in C\) 。（存在 c 使 \(cx_{m} + y \succ_{C} 0\) ；利用前一习题的 (e) 部分得出结论。）
\end{enumerate}

\(d)\) 设 \(\mathrm{C}^{\circ}\subset\mathrm{E}^{\prime}\) 是 C 的极集。对每个 \(x\in\mathring{\mathrm{C}}\) 与每个 \(\xi\in\mathrm{C}^{\circ}-\{0\}\) ，有 \(\langle x,\xi\rangle>0\) 。

\begin{enumerate}
\def\labelenumi{\alph{enumi})}
\setcounter{enumi}{4}
\tightlist
\item
  C 的极集 C° 在 \(^{t}u\colon E'\to E'\) 下不变；对每个 \(\lambda\in R_{+}\) （满足 \(\lambda<\lambda_{m}\) ），自同态 \(1_{E'}-\lambda^{t}u\) （\(E'\) 的）是一个同构。
\end{enumerate}

\begin{enumerate}
\def\labelenumi{\arabic{enumi})}
\setcounter{enumi}{29}
\tightlist
\item
  我们沿用前一习题的假设，即 u 是线性、紧的，且假设条件 (i) 至 (v) 成立。自同态 \(^{t}u\) 也是紧的。
\end{enumerate}

\begin{enumerate}
\def\labelenumi{\alph{enumi})}
\item
  设 \(\eta \in \mathrm{C}^{\circ} - \{0\}\) 。对每个 \(\lambda \in \mathbf{R}_{+}\) （满足 \(\lambda < \lambda_{\mathrm{m}}\) ），唯一解 \(\xi(\lambda)\) （方程 \(\xi - \lambda^{t}u(\xi) = \eta\) 的）属于 \(\mathrm{C}^{\circ}\) 。（若 \(x \in \mathrm{C}\) ，设 \(y \in \mathrm{C}\) 满足 \(y - \lambda u(y) = x\) ；则 \(\langle y, \eta \rangle = \langle x, \xi \rangle\) 。）
\item
  设 \(x_{0} \in \mathring{C}\) 。映射 \(\lambda \mapsto \langle x_{0}, \xi(\lambda) \rangle\) 当 \(\lambda\) 趋于 \(\lambda_{m}\) 时趋于无穷。证明值 \(\xi(\lambda)/\langle x_{0}, \xi(\lambda) \rangle\) 所成的集有一个聚点 \(\xi_{m}\) 满足：
\end{enumerate}

\begin{enumerate}
\def\labelenumi{(\roman{enumi})}
\item
  \(\xi_{\mathrm{m}}\in \mathrm{C}^{\circ}\) ;
\item
  \(\langle x_{0},\xi_{m}\rangle=1;\)
\item
  \(\xi_{\mathrm{m}} = \lambda_{\mathrm{m}}^{t}u(\xi_{\mathrm{m}})\) .
\end{enumerate}

\begin{enumerate}
\def\labelenumi{\alph{enumi})}
\setcounter{enumi}{2}
\item
  设 \(f \in E\) 满足 \(\langle f, \xi_{m} \rangle = 0\) 。对 \(\lambda \in R_{+}\) （满足 \(\lambda < \lambda_{m}\) ），用 \(x(\lambda)\) 表示唯一解 \(x\) （方程 \(x - \lambda u(x) = f\) 在 E 中的）。证明 \(x(\lambda)\) 所成的集在 E 中有界（设 d 是定义 E 的拓扑的距离；若存在序列 \(\lambda_{n} \to \lambda_{m}\) 使 \(x(\lambda_{n})\) 趋于无穷，置 \(\hat{x}_{n} = x(\lambda_{n}) / d(0, x(\lambda_{n}))\) ，并证明序列 \((\hat{x}_{n})\) 收敛于一个元素 \(\hat{x}\) ，它满足 \(d(0, \hat{x}) = 1\) 、\(\hat{x} = \lambda u(\hat{x})\) 与 \(\langle \hat{x}, \xi_{m} \rangle = 0\) ，这与必有 \(\hat{x} \in \mathring{C}\) 的事实矛盾）。
\item
  存在唯一的元素 \(x \in E\) 满足 \(f = x - \lambda_{\mathrm{m}} u(x)\) 与 \(\langle x, \xi_{\mathrm{m}} \rangle = 0\) 。（关于存在性，利用前一问题证明存在一个序列 \((\lambda_{n})\) 收敛于 \(\lambda_{m}\) ，使 \((x(\lambda_{n}))\) 收敛于一个满足所需条件的元素 x。）
\item
  自同态 \(1_{\mathrm{E}'} - \lambda_{\mathrm{m}}{}^{t}u\) （\(\mathrm{E}'\) 的）通过过渡到子空间诱导极集 \(x_{\mathrm{m}}^{\circ}\) （\(x_{\mathrm{m}}\) 的）上的一个双射。
\item
  特征空间 \(\operatorname{Ker}(1_{\mathrm{E}} - \lambda_{\mathrm{m}}u)\) 的维数为 1，且等于 \(1_{E} - \lambda_{m}u\) 的像的极集。
\end{enumerate}

第四章

